\subsection{Application: homology and cohomology of groups}

Let G be a group, \(\mathbf{Z}^{(G)}\) its algebra over Z (III, p.~19). Recall (cf.~III, p.~20, example) that if M is a commutative group, it is equivalent to giving an action of G on M (that is, a homomorphism \(\tau: G \to \text{Aut}(M)\) ), or a structure of \(\mathbf{Z}^{(G)}\) left module on the additive group M. In particular, one will consider the group Z as a \(\mathbf{Z}^{(G)}\) -module on the left, endowing it with the trivial action.

DEFINITION 1. --- Let M be a \(\mathbf{Z}^{(G)}\) left (resp. right) module, \(n\) be an integer \(\geqslant 0\) . The group \(\operatorname{Ext}_{\mathbf{Z}^{(G)}}^{n}(\mathbf{Z}, \mathbf{M})\) (resp. \(\operatorname{Tor}_{n}^{\mathbf{Z}^{(G)}}(\mathbf{M}, \mathbf{Z})\) ) is denoted \(\mathrm{H}^{n}(\mathrm{G}, \mathbf{M})\) (resp. \(\mathrm{H}_{n}(\mathrm{G}, \mathbf{M}))\) and called \(n\) -th cohomology group (resp. homology group) of \(\mathbf{G}\) with coefficients in \(\mathbf{M}\) .

The standard resolution (X, p.~58) B(Z \(^{(G)}\) , Z) is a free resolution of the Z \(^{(G)}\) -module Z; it follows that the groups H \(^{n}\) (G, M) (resp. H \(_{n}\) (G, M)) are identified with the homology groups of the complex:

\[
\operatorname{Hom} _ {\mathbf {Z} ^ {(G)}} \left(\mathrm{B} (\mathbf {Z} ^ {(G)}, \mathbf {Z}), \mathrm{M}\right) \quad \left(\text { resp. } \mathrm{M} \otimes_ {\mathbf {Z} ^ {(G)}} \mathrm{B} (\mathbf {Z} ^ {(G)}, \mathbf {Z})\right).
\]

Using the canonical isomorphism of \((\mathbf{Z}^{(\mathbf{G})})^{\otimes n}\) onto \(\mathbf{Z}^{(\mathbf{G}^{n})}\) (III, p.~36) and the properties of extension of scalars (II, p.~82), we conclude that \(\mathrm{H}^n (\mathrm{G},\mathrm{M})\) is canonically isomorphic to the homology group of ascending degree \(n\) of the complex \(\mathrm{C(G,M)}\) defined in the following way: \(\mathbf{C}^n (\mathbf{G},\mathbf{M}) = 0\) for \(n < 0\) ; for \(n\geqslant 0\) , \(\mathbf{C}^n (\mathbf{G},\mathbf{M})\) is the \(\mathbf{Z}\) -module of maps from \(\mathbf{G}^n\) to \(\mathbf{M}\) ; for \(n\geqslant 0\) , the differential \(d^n:\mathbf{C}^n (\mathbf{G},\mathbf{M})\to \mathbf{C}^{n + 1}(\mathbf{G},\mathbf{M})\) is given by

\[
(d ^ {n} f) (g _ {0}, \dots , g _ {n}) = g _ {0} \cdot f (g _ {1}, \dots , g _ {n}) + \sum_ {i = 0} ^ {n - 1} (- 1) ^ {i + 1} f (g _ {0}, \dots , g _ {i} g _ {i + 1}, \dots , g _ {n})
\]

\[
+ (- 1) ^ {n + 1} f \left(g _ {0}, \dots , g _ {n - 1}\right)
\]

whatever \(f\) in \(\mathbf{C}^{n}(\mathbf{G},\mathbf{M})\) and \(g_{0},\ldots,g_{n}\) in \(\mathbf{G}\) .

Similarly, \(\mathrm{H}_{n}(\mathrm{G},\mathrm{M})\) is identified with the homology group of degree n of the complex \(\mathrm{C}'(\mathrm{G},\mathrm{M})\) , where \(\mathrm{C}_{n}^{\prime}(\mathrm{G},\mathrm{M})=\mathrm{M}\otimes_{\mathbf{Z}}\mathbf{Z}^{(\mathbf{G}^{n})}\) for \(n\geqslant0\) , \(\mathrm{C}_{n}^{\prime}(\mathrm{G},\mathrm{M})=0\) for n\textless0, the differential \(d_{n}:C_{n}^{\prime}(\mathrm{G},\mathrm{M})\to C_{n-1}^{\prime}(\mathrm{G},\mathrm{M})\) being defined by:

\[
d _ {n} (m \otimes e _ {g _ {1}, \dots , g _ {n}}) = m. g _ {1} \otimes e _ {g _ {2}, \dots , g _ {n}} + \sum_ {i = 1} ^ {n - 1} (- 1) ^ {i} m \otimes e _ {g _ {1}, \dots , g _ {i} g _ {i + 1}, \dots , g _ {n}}
\]

\[
+ (- 1) ^ {n} m \otimes e _ {g _ {1}, \dots , g _ {n - 1}}
\]

whatever \(n \geqslant 1\) , \(m\) in M and \(g_1, \ldots, g_n\) in G.

Examples. --- 1) It follows directly from the definition that H \(^{0}\) (G, M) is isomorphic to the submodule of elements of M invariant under the action of G, and H \(_{0}\) (G, M) to the quotient module of M by the submodule generated by the elements m.g - m for m ∈ M, g ∈ G.

\begin{enumerate}
\def\labelenumi{\arabic{enumi})}
\setcounter{enumi}{1}
\tightlist
\item
  It follows from the preceding that \(H^{1}(G, M)\) is isomorphic to the Z-module \(Z^{1}(G, M)/B^{1}(G, M)\) , where \(Z^{1}(G, M)\) is the Z-module of functions f from G to M satisfying:
\end{enumerate}

\[
f \left(g _ {1}, g _ {2}\right) = g _ {1} \cdot f \left(g _ {2}\right) + f \left(g _ {1}\right) \quad \text { for   all } \quad g _ {1}, g _ {2} \text { in } G,
\]

and \(\mathbf{B}^{1}(\mathbf{G},\mathbf{M})\) is the sub \(\mathbf{Z}\) -module of \(\mathbf{Z}^{1}(\mathbf{G},\mathbf{M})\) consisting of \(f\) for which there exists an element \(m\) of \(\mathbf{M}\) such that:

\[
f (g) = g. m - m \quad \text { for   every } \quad g \in G.
\]

One sometimes says that \(Z^{1}(G, M)\) is the Z-module of crossed homomorphisms from G to M, and \(B^{1}(G, M)\) the submodule of principal crossed homomorphisms.

Let us denote by \(\tau: \mathbf{G} \to \operatorname{Aut}(\mathbf{M})\) the homomorphism deduced from the action of \(\mathbf{G}\) ; consider the external semidirect product \(\mathbf{M} \times_{\tau} \mathbf{G}\) and the extension \(\xi_{\tau}: \mathbf{M} \xrightarrow{i} \mathbf{M} \times_{\tau} \mathbf{G} \xrightarrow{p} \mathbf{G}\) (I, p.~64). Let \(e: \mathbf{G} \to \mathbf{M} \times_{\tau} \mathbf{G}\) be a map such that \(p \circ e = 1_{\mathbf{G}}\) ; we have \(e = (f, 1_{\mathbf{G}})\) , where \(f \in C^{1}(\mathbf{G}, \mathbf{M})\) . For \(e\) to be a homomorphism (that is, a section of the extension \(\xi_{\tau}\) ), it is necessary and sufficient that \(f \in Z^{1}(\mathbf{G}, \mathbf{M})\) . For two sections of \(\xi_{\tau}\) to be conjugate by an element of \(i(\mathbf{M})\) , it is necessary and sufficient that the corresponding crossed homomorphisms have the same class in \(H^{1}(\mathbf{G}, \mathbf{M})\) .

When G acts trivially on M, we have \(\mathbf{B}^{1}(\mathbf{G},\mathbf{M})=0\) and \(\mathbf{H}^{1}(\mathbf{G},\mathbf{M})\) is isomorphic to the Z-module of group homomorphisms from G to M.

\begin{enumerate}
\def\labelenumi{\arabic{enumi})}
\setcounter{enumi}{2}
\tightlist
\item
  Similarly \(H^{2}(G, M)\) is isomorphic to the Z-module \(Z^{2}(G, M)/B^{2}(G, M)\) , where \(Z^{2}(G, M)\) is the Z-module of maps f from G × G to M, satisfying:
\end{enumerate}

\[
g _ {1} \cdot f (g _ {2}, g _ {3}) - f (g _ {1} g _ {2}, g _ {3}) + f (g _ {1}, g _ {2} g _ {3}) - f (g _ {1}, g _ {2}) = 0
\]

whatever \(g_{1}, g_{2}, g_{3}\) in \(\mathbf{G}\) , and \(\mathbf{B}^{2}(\mathbf{G}, \mathbf{M})\) is the sub- \(\mathbf{Z}\) -module of \(\mathbf{Z}^{2}(\mathbf{G}, \mathbf{M})\) consisting of \(f\) for which there exists a mapping \(h\) of \(\mathbf{G}\) to \(\mathbf{M}\) such that:

\[
f (g _ {1}, g _ {2}) = g _ {1}. h (g _ {2}) - h (g _ {1} g _ {2}) + h (g _ {1})
\]

whatever \(g_{1}, g_{2}\) in G.

We thus recover the definition of the group \(H^{2}(G, M)\) given in VIII, App. It follows in particular that there exists a canonical isomorphism of \(H^{2}(G, M)\) on the group of extension classes of G by M (loc. cit.).

\begin{enumerate}
\def\labelenumi{\arabic{enumi})}
\setcounter{enumi}{3}
\tightlist
\item
  Let M be a Z-module, which we consider as a \(\mathbf{Z}^{(\mathrm{G})}\) -module on the right by letting G act trivially. The group \(\mathrm{H}_{1}(\mathrm{G},\mathrm{M})\) is isomorphic to the quotient of \(\mathrm{M}\otimes_{\mathbf{Z}}\mathbf{Z}^{(\mathrm{G})}\) by the sub-Z-module generated by the elements \(m\otimes(e_{g_{1}g_{2}}-e_{g_{1}}-e_{g_{2}})\) for m in M, \(g_{1}, g_{2}\) in G; it follows easily that \(\mathrm{H}_{1}(\mathrm{G},\mathrm{M})\) is isomorphic to \(\mathrm{M}\otimes_{\mathbf{Z}}(\mathrm{G}/(\mathrm{G},\mathrm{G}))\) .
\end{enumerate}

Let us denote by \(\sigma\) the anti-automorphism of \(\mathbf{Z}^{(G)}\) defined by \(\sigma(e_g) = e_{g-1}\) for \(g \in G\) . Every \(\mathbf{Z}^{(G)}\) -module on the left can be considered as a \(\mathbf{Z}^{(G)}\) -module on the right by means of \(\sigma\) , and conversely. This allows us, for example, to define the groups \(\mathrm{H}_q(\mathrm{G}, \mathrm{M})\) for a \(\mathbf{Z}^{(G)}\) -module on the left M, by setting \(\mathrm{H}_q(\mathrm{G}, \mathrm{M}) = \mathrm{H}_q(\mathrm{G}, \sigma_*(\mathrm{M})) = \mathrm{Tor}_q^{\mathbf{Z}^{(G)}}(\mathbf{Z}, \mathbf{M})\) .

Lemma 1. --- Let M be a Z-module; denote by M \(^{G}\) the group Hom \(_{Z}\) (Z \(^{(G)}\) , M) equipped with its natural Z \(^{(G)}\) -module structure on the left. Then:

\[
\mathrm{H} ^ {i} (\mathrm{G}, \mathrm{M} ^ {\mathrm{G}}) = 0 \quad \text { for } \quad i \geqslant 1.
\]

It follows in fact from prop. 8, b) (X, p.~109), applied with A = N = Z \(^{(G)}\) and B = Q = Z, that we have a canonical isomorphism:

\[
\operatorname{Ext} _ {\mathbf {Z} ^ {(G)}} \left(\mathbf {Z}, \mathrm{M} ^ {\mathrm{G}}\right)\rightarrow \operatorname{Ext} _ {\mathbf {Z}} (\mathbf {Z}, \mathrm{M})
\]

whence the lemma.

PROPOSITION 11. --- Let L be a finite Galois extension of a commutative field K, with Galois group G.

\begin{enumerate}
\def\labelenumi{\alph{enumi})}
\item
  One has \(\mathrm{H}^i (\mathbf{G},\mathbf{L}) = 0\) for \(i\geqslant 1\)
\item
  One has \(\mathrm{H}^1 (\mathbf{G},\mathbf{L}^*) = 0\)
\item
  The group \(\mathrm{H}^2 (\mathbf{G},\mathbf{L}^*)\) is canonically isomorphic to the group Br(K, L) (VIII, § 13).
\end{enumerate}

The normal basis theorem (V, §10, no. 9, th. 6) shows that L is isomorphic as \(\mathbf{Z}^{(G)}\) -module to \(\mathbf{K}^{\mathrm{G}} = \operatorname{Hom}_{\mathbf{Z}}(\mathbf{Z}^{(G)}, \mathbf{K})\) ; the assertion \(a)\) then follows from Lemma 1. In view of Example 2, the assertion \(b)\) follows from V, §10, no. 5, cor. 1 to prop. 9; finally, the assertion \(c)\) was proved in VIII, §13.

